%%% chapter3.tex — Глава 3
\section{Программная реализация системы \texttt{NetStego}}

В данной главе представлена программная реализация системы \texttt{NetStego}. Полные листинги исходного кода приведены в Приложениях~А--В.

\subsection{Используемые инструменты разработки}

В качестве языка программирования выбран Python~3.11, что обусловлено наличием зрелых библиотек для работы с сетевыми протоколами, криптографией и машинным обучением. Основные компоненты технологического стека представлены в таблице~\ref{tab:tech_stack}.

\begin{table}[htb]
\centering
\caption{Технологический стек системы \texttt{NetStego}}
\label{tab:tech_stack}
\small
\begin{tabular}{|l|l|p{7cm}|}
\hline
Библиотека & Версия & Назначение \\
\hline
Scapy~[21]          & 2.5+  & Создание, отправка и анализ пакетов \\
Cryptography~[22]   & 42+   & AES-256-GCM шифрование \\
argon2-cffi          & 23+   & Деривация ключа (Argon2id, RFC~9106~[9]) \\
Scikit-learn~[23]    & 1.4+  & Isolation Forest, Random Forest (ML) \\
SciPy                & 1.12+ & Статистические тесты ($\chi^2$, KS) \\
Click                & 8.1+  & Интерфейс командной строки \\
SQLAlchemy           & 2.0+  & ORM для SQLite базы статистики \\
Matplotlib           & 3.8+  & Генерация графиков и диаграмм \\
Rich                 & 13+   & Форматирование табличного вывода \\
Pydantic             & 2.0+  & Валидация конфигурации \\
\hline
\end{tabular}
\end{table}

Система поддерживает работу на платформах Windows (с установленным драйвером Npcap) и Linux (с привилегиями CAP\_NET\_RAW для захвата и отправки пакетов).

Структура проекта организована в виде Python-пакета \texttt{netstego} со следующей модульной организацией:

\begin{enumerate}[label=\arabic*)]
    \item \texttt{core/} --- криптографическое ядро (шифрование, фрагментация, сборка);
    \item \texttt{channels/} --- реализации стеганографических каналов;
    \item \texttt{network/} --- модули отправки и приёма данных с интеграцией сбора статистики;
    \item \texttt{detection/} --- модуль обнаружения скрытых каналов и анализа вредоносных вложений;
    \item \texttt{stats/} --- сбор статистики, база данных, отчёты и генерация графиков;
    \item \texttt{benchmarks.py} --- бенчмарки конвейера и мониторинг точности обнаружения;
    \item \texttt{cli.py} --- точка входа CLI;
    \item \texttt{config.py} --- модели конфигурации.
\end{enumerate}

\subsection{Применённые подходы при реализации системы}

\subsubsection{Архитектура конвейера обработки данных}

Архитектура системы построена в виде конвейера (pipeline), последовательно обрабатывающего данные на стороне отправителя и приёмника. Конвейер отправки включает этапы:

\begin{enumerate}[label=\arabic*)]
    \item загрузка исходного файла и 256-битного ключа шифрования;
    \item шифрование файла алгоритмом AES-256-GCM с генерацией уникального 12-байтного nonce;
    \item фрагментация шифротекста на чанки с добавлением протокольного заголовка NST (17~байт);
    \item фрейминг --- добавление 4-байтового префикса длины к каждому чанку;
    \item кодирование потока байт в сетевые пакеты выбранным стеганографическим каналом;
    \item отправка пакетов с заданной скоростью (PPS) и джиттером;
    \item запись метаданных сессии и пакетов в базу данных статистики.
\end{enumerate}

Конвейер приёма выполняет обратные преобразования: захват пакетов (live или из PCAP), декодирование каналом, разбор фреймов, сборка чанков с обработкой out-of-order, расшифрование, запись файла. На каждом этапе ведётся протоколирование в базе статистики с фиксацией статуса сессии (completed/failed).

Реализация конвейера отправки представлена в Приложении~А.

\subsubsection{Криптографическая защита}

Для обеспечения конфиденциальности и целостности передаваемых данных применяется алгоритм аутентифицированного шифрования AES-256-GCM, описанный в NIST~SP~800-38D~[8]. Выбор данного алгоритма обусловлен следующими свойствами:

\begin{enumerate}[label=\arabic*)]
    \item обеспечение одновременно конфиденциальности и аутентификации (AEAD);
    \item 256-битный ключ, обеспечивающий устойчивость к атакам перебора;
    \item 128-битный тег аутентификации, гарантирующий обнаружение модификации;
    \item 96-битный nonce (12~байт), генерируемый криптографически стойким генератором случайных чисел.
\end{enumerate}

Формат зашифрованных данных:
\begin{center}
\texttt{nonce (12 байт) || ciphertext || tag (16 байт)}
\end{center}

Генерация ключа осуществляется функцией \texttt{secrets.token\_bytes(32)}, обеспечивающей криптографически стойкую случайность. Для деривации ключа из пароля используется функция Argon2id (RFC~9106~[9]) с параметрами: time\_cost~=~3, memory\_cost~=~65536~КиБ, что обеспечивает устойчивость к атакам по словарю и на специализированном оборудовании.

Реализация модуля шифрования приведена в Приложении~Б.

\subsubsection{Протокол фрагментации и сборки}

Зашифрованные данные разбиваются на чанки фиксированного размера, каждый из которых снабжается внутренним протокольным заголовком. Формат заголовка чанка (17~байт) представлен в таблице~\ref{tab:chunk_header}.

\begin{table}[htb]
\centering
\caption{Формат заголовка чанка протокола NST}
\label{tab:chunk_header}
\begin{tabular}{|c|c|l|}
\hline
Смещение & Размер, байт & Поле \\
\hline
0  & 4 & Магическая последовательность \texttt{NST\textbackslash x00} \\
4  & 4 & Порядковый номер (SEQ, uint32 big-endian) \\
8  & 4 & Общее количество чанков (TOTAL, uint32 BE) \\
12 & 4 & Контрольная сумма CRC32 данных (uint32 BE) \\
16 & 1 & Флаги (FLAGS, uint8) \\
17 & $N$ & Данные чанка \\
\hline
\end{tabular}
\end{table}

Флаги кодируются побитово: 0x01 --- первый чанк (данные начинаются с SHA-256 хеша исходного файла, 32~байта), 0x02 --- последний чанк, 0x04 --- запрос повторной передачи (NACK).

Включение SHA-256 хеша в первый чанк обеспечивает сквозную верификацию целостности файла после сборки. CRC32 каждого чанка позволяет выявить повреждение данных на уровне отдельных фрагментов.

Реализация функции фрагментации приведена в Приложении~Б.

Класс \texttt{Reassembler} реализует сборку чанков с поддержкой приёма не по порядку (out-of-order). Чанки хранятся в словаре с ключом --- порядковым номером. При каждом добавлении проверяется CRC32, после получения всех чанков выполняется сортировка по номерам, извлечение SHA-256 из первого чанка и верификация целостности восстановленных данных.

\subsubsection{Абстракция каналов сокрытия}

Все каналы реализованы как подклассы абстрактного базового класса \texttt{StegoChannel}, определяющего единый интерфейс. Класс включает четыре обязательных метода:

\begin{enumerate}[label=\arabic*)]
    \item \texttt{encode(data, dest\_ip)} --- кодирование данных в список пакетов Scapy;
    \item \texttt{decode(packets)} --- извлечение скрытых данных из списка пакетов;
    \item \texttt{bits\_per\_packet} --- количество бит данных на один пакет;
    \item \texttt{protocol\_filter} --- BPF-фильтр для перехвата пакетов.
\end{enumerate}

Такая архитектура обеспечивает расширяемость: для добавления нового канала достаточно реализовать подкласс \texttt{StegoChannel} с четырьмя методами.

\subsection{Реализация функциональных возможностей}

Система \texttt{NetStego} предоставляет интерфейс командной строки (CLI) с семью командами: \texttt{keygen} --- генерация 256-битного ключа, \texttt{send} --- шифрование, фрагментация и отправка данных, \texttt{receive} --- приём и восстановление данных, \texttt{detect} --- анализ трафика на наличие скрытых каналов и вредоносных вложений, \texttt{stats} --- просмотр статистики сессий, \texttt{benchmark} --- бенчмарк канала, \texttt{report} --- полный отчёт с генерацией графиков.

\subsubsection{Передача данных}

Команда \texttt{send} поддерживает два режима ввода: передача файла (\texttt{-{}-file}) и передача текстовой строки (\texttt{-{}-text}), что позволяет отправлять произвольные типы данных --- текстовые сообщения, изображения, документы, криптографические ключи. Параметр \texttt{-{}-redundancy~N} активирует механизм устойчивости к потере пакетов: каждый чанк дублируется $N$~раз, а дубликаты перемешиваются для равномерного распределения по потоку. На стороне приёмника класс \texttt{Reassembler} автоматически дедуплицирует чанки по порядковому номеру. Параметр \texttt{-{}-pcap-out} сохраняет отправленные пакеты в формат PCAP. Параметр \texttt{-{}-db} позволяет указать путь к базе данных статистики для протоколирования сессии.

Конвейер отправки определяет размер чанка в зависимости от ёмкости канала: для высокоёмких каналов (ICMP, 1472~байта) чанк помещается в один пакет, для низкоёмких (IP~ID, 2~байта) используется фиксированный размер 128~байт. Регулировка скорости реализована с помощью межпакетных задержек с джиттером ($\pm10$\,\%) для повышения скрытности.

Конвейер приёма выполняет обратные операции: захват пакетов с BPF-фильтром, декодирование каналом, извлечение и сборка чанков с верификацией CRC32 и SHA-256, расшифрование AES-256-GCM. На каждом этапе, где возможна ошибка (отсутствие пакетов, ошибка декодирования, неполная сборка, ошибка расшифровки), сессия помечается статусом <<failed>> в базе статистики с фиксацией количества обработанных пакетов.

\subsubsection{Стеганографические каналы}

Каждый канал реализован как подкласс абстрактного базового класса \texttt{StegoChannel} с методами \texttt{encode}/\texttt{decode}. Характеристики каналов представлены в таблице~\ref{tab:channels}.

\begin{table}[htb]
\centering
\caption{Характеристики стеганографических каналов}
\label{tab:channels}
\small
\begin{tabular}{|l|l|c|l|}
\hline
Канал & Поле-носитель & Бит/пак. & Механизм кодирования \\
\hline
ICMP Payload   & Payload ICMP Echo & 11\,776 & Length prefix (2~Б) \\
IP ID          & IPv4 Identification & 16 & Прямая запись uint16 \\
TCP ISN        & TCP Seq Number & 32 & SYN с инкр. dport \\
DNS QNAME      & Поддомен DNS & 240 & Base32-кодирование \\
TCP Timestamp  & TCP TSval Option & 32 & Прямая запись uint32 \\
\hline
\end{tabular}
\end{table}

Канал ICMP Payload помещает данные в поле полезной нагрузки ICMP Echo Request (до 1470~байт на пакет с учётом length prefix). Канал IP~ID встраивает 2~байта в поле \texttt{IP.id}, используя терминирующий пакет с \texttt{id=0xFFFF} для передачи общей длины. Канал TCP~ISN --- 4~байта в поле \texttt{TCP.seq} SYN-пакетов. Канал TCP~Timestamp встраивает 4~байта в поле TSval TCP Timestamp Option.

Канал DNS~QNAME кодирует данные в base32 в субдоменах DNS-запросов (до 30~байт на пакет). Доменный суффикс является конфигурируемым параметром (по умолчанию --- реалистичное доменное имя вместо фиксированной сигнатуры), что снижает вероятность обнаружения сигнатурным методом. Терминирующий DNS-запрос содержит маркер <<term.>> с закодированной общей длиной данных.

\subsubsection{Модуль обнаружения скрытых каналов}

Модуль обнаружения объединяет четыре метода анализа: статистический, машинное обучение (unsupervised и supervised) и сигнатурный с анализом содержимого payload.

\textbf{Статистический анализ} включает девять тестов: критерий $\chi^2$ для оценки равномерности распределения значений полей, тест Колмогорова--Смирнова для проверки экспоненциальности межпакетных задержек, нормализованная энтропия Шеннона для выявления случайных данных в полях заголовков, коэффициент вариации IPD для обнаружения регулярных задержек, проверка последовательности IP~ID (в нормальном Windows-трафике ID инкрементируется на 1--3), проверка монотонности TSval, анализ энтропии DNS QNAME, анализ размера ICMP payload и анализ побайтовой энтропии payload (зашифрованные данные имеют $H > 7{,}0$~бит/байт).

\textbf{ML-обнаружение (unsupervised)} реализовано на основе алгоритма Isolation Forest с 9~признаками: значение поля, дельта от предыдущего значения, межпакетная задержка, размер пакета, скользящее среднее и стандартное отклонение значений поля (окно~5), скользящее среднее IPD, энтропия payload (бит/байт) и длина payload. Стандартизация признаков выполняется с помощью StandardScaler. Аномалия фиксируется при доле аномальных пакетов, превышающей $2 \cdot \texttt{contamination}$.

\textbf{ML-обнаружение (supervised)} реализовано на основе алгоритма Random Forest (100~деревьев, max\_depth~=~10) с теми же 9~признаками. Классификатор обучается на размеченных данных: нормальный трафик (класс~0) и стеганографический трафик (класс~1). Данный подход используется в модуле мониторинга для оценки теоретически достижимой точности классификации каналов.

\textbf{Сигнатурный анализ} выполняет поиск в полезной нагрузке пакетов магической последовательности \texttt{NST\textbackslash x00} протокола фрагментации (уверенность 100\,\%), паттерна <<length prefix + magic>> (уверенность 95\,\%), base32-паттернов длиной более 20~символов в DNS QNAME (уверенность 85\,\%) и терминационных маркеров DNS (уверенность 90\,\%).

\textbf{Анализ содержимого payload (обнаружение вредоносных вложений).} Дополнительно к обнаружению факта стеганографии система анализирует содержимое извлечённых данных на наличие индикаторов вредоносных вложений. Данный модуль сканирует payload пакетов на наличие 20 сигнатур, сгруппированных по категориям (таблица~\ref{tab:content_sigs}).

\begin{table}[htb]
\centering
\caption{Сигнатуры анализа содержимого payload}
\label{tab:content_sigs}
\begin{tabular}{|l|l|c|}
\hline
Категория & Сигнатуры & Уверенность \\
\hline
Исполняемые файлы  & MZ (PE/Windows), \texttt{\textbackslash x7fELF} (Linux) & 0{,}90--0{,}95 \\
Скрипты            & Shebang (\texttt{\#!/bin/}), PowerShell, cmd.exe & 0{,}75--0{,}80 \\
Веб-атаки          & \texttt{<script>}, \texttt{eval()}, \texttt{exec()} & 0{,}70 \\
Шеллкод            & NOP sled (\texttt{\textbackslash x90}$\times$8), INT3 sled & 0{,}80--0{,}85 \\
Архивы             & ZIP (PK), GZIP, RAR & 0{,}50--0{,}60 \\
Макро-документы    & OLE2 (Office), PDF & 0{,}50--0{,}70 \\
Обфускация         & base64.b64decode & 0{,}60 \\
\hline
\end{tabular}
\end{table}

Данный подход позволяет не только выявить факт использования скрытого канала, но и классифицировать тип скрытых данных, что существенно для оценки уровня угрозы: передача исполняемого файла или шеллкода представляет значительно более высокий риск, чем передача текстового сообщения.

Оркестратор вычисляет итоговую оценку подозрительности по взвешенной формуле:
\begin{equation}
C = 0{,}4 \cdot I_{\text{stat}} + 0{,}3 \cdot I_{\text{ml}} + 0{,}3 \cdot I_{\text{sig}},
\label{eq:confidence}
\end{equation}
где $I_{\text{stat}} = 1$, если более 50\,\% статистических тестов выявили аномалию или обнаружен хотя бы один сильный индикатор (аномалия TSval, DNS QNAME, размер ICMP payload, энтропия payload, последовательность IP~ID); $I_{\text{ml}} = 1$, если доля аномалий по Isolation Forest превышает $2 \cdot \texttt{contamination}$; $I_{\text{sig}} = 1$, если обнаружены сигнатуры протокола или вредоносного содержимого. Трафик считается подозрительным при $C \geq 0{,}4$. Реализация модуля обнаружения приведена в Приложении~В.

\subsubsection{Система сбора статистики и визуализации}

Модуль \texttt{StatsCollector} интегрирован в конвейеры отправки и приёма и обеспечивает протоколирование каждой сессии передачи данных. При начале операции создаётся запись сессии с метаданными (направление, канал, IP-адреса, размер файла), в ходе обработки фиксируется информация о каждом пакете (порядковый номер, имя поля, размер), по завершении сессия помечается итоговым статусом.

База данных SQLite содержит три таблицы: \texttt{sessions} (метаданные сессий), \texttt{packets\_log} (журнал пакетов), \texttt{anomaly\_scores} (результаты обнаружения). Поддерживается экспорт отчётов в форматах JSON, CSV и HTML.

Модуль визуализации (\texttt{stats/charts.py}) использует библиотеку Matplotlib (backend Agg для работы без GUI) и генерирует шесть типов графиков:

\begin{enumerate}[label=\arabic*)]
    \item \textbf{Сравнение пропускной способности} --- столбчатая диаграмма (лог.~шкала) с подписями значений в MB/s или KB/s;
    \item \textbf{Декомпозиция времени} --- горизонтальная составная диаграмма по этапам конвейера (шифрование, фрагментация, кодирование, декодирование, сборка, расшифрование);
    \item \textbf{Пакеты vs размер данных} --- лог-лог график зависимости числа пакетов от размера payload;
    \item \textbf{Пропускная способность vs размер данных} --- график масштабирования throughput;
    \item \textbf{Точность обнаружения} --- диаграмма Precision/Recall/F1 по каналам;
    \item \textbf{Устойчивость к потерям} --- диаграмма восстановления при потерях пакетов.
\end{enumerate}

Команда CLI \texttt{report} запускает полный цикл бенчмарков (5~каналов $\times$ N~размеров $\times$ K~повторов), опционально --- мониторинг обнаружения, сохраняет результаты в JSON (\texttt{results/data/}) и генерирует графики в PNG (\texttt{results/charts/}).

\subsection{Тестирование системы}

\subsubsection{Среда тестирования}

Тестирование проводилось в следующей среде: ОС~--- Windows~10~Pro и Ubuntu~22.04~LTS, сеть~--- loopback-интерфейс (127.0.0.1), фреймворк тестирования~--- pytest с плагинами для оценки покрытия. Для визуального контроля корректности формирования пакетов использовался анализатор Wireshark.

\subsubsection{Модульные тесты}

Модульные тесты (unit tests) проверяют корректность работы отдельных компонентов системы. Реализовано 134 модульных теста, покрывающих следующие области:

\begin{enumerate}[label=\arabic*)]
    \item \textbf{Криптографический модуль} (\texttt{test\_crypto.py}, 12~тестов): корректность шифрования/дешифрования, обнаружение повреждённых данных (InvalidTag), проверка формата (nonce~+ ciphertext~+ tag), деривация ключа Argon2id.
    \item \textbf{Фрагментация и сборка} (\texttt{test\_fragmentation.py}, 12~тестов; \texttt{test\_pipeline.py}, 17~тестов): разбиение на чанки, CRC32, сборка в произвольном порядке, обнаружение повреждённых чанков, полный pipeline для всех каналов.
    \item \textbf{Каналы сокрытия} (5 файлов тестов, 40 тестов): roundtrip-тестирование каждого канала по формуле $\texttt{decode}(\texttt{encode}(data)) = data$ для данных различного объёма, проверка конфигурируемого суффикса DNS.
    \item \textbf{Модуль обнаружения} (\texttt{test\_detection.py}, 33 теста): корректность вычисления энтропии, $\chi^2$-теста, работы Isolation Forest и Random Forest на синтетических данных, анализ содержимого payload (PE, ELF, скрипты, шеллкод, архивы, чистые данные).
    \item \textbf{Бенчмарки и мониторинг} (\texttt{test\_benchmarks.py}, 35 тестов): pipeline-бенчмарки для всех каналов, тесты избыточности, мониторинг обнаружения (включая supervised ML), генераторы нормального и стеганографического трафика, метрики DetectionMetrics.
    \item \textbf{Графики} (\texttt{test\_charts.py}, 17 тестов): генерация каждого из 6 типов графиков, проверка формата PNG и размера файла.
    \item \textbf{Модуль статистики} (\texttt{test\_stats.py}, 5 тестов): создание и чтение записей в SQLite, экспорт отчётов.
\end{enumerate}

\subsubsection{Интеграционные тесты}

Интеграционные тесты (68~тестов) проверяют работу полного конвейера:

\begin{enumerate}[label=\arabic*)]
    \item \textbf{Целостность данных} (\texttt{test\_data\_integrity.py}, 67 тестов): полный roundtrip через PCAP для всех 5~каналов с различными типами данных (ASCII, UTF-8, бинарные), размерами (1~байт -- 100~КБ), побайтовое сравнение SHA-256 хешей.
    \item \textbf{PCAP-конвейер} (\texttt{test\_pcap\_pipeline.py}, 15 тестов): отправка с сохранением в PCAP, приём из PCAP для каждого канала.
    \item \textbf{CLI-интерфейс} (\texttt{test\_cli.py}, 10 тестов): тестирование команд CLI через \texttt{CliRunner}, проверка параметров и кодов возврата.
    \item \textbf{Статистика} (\texttt{test\_stats\_integration.py}, 5 тестов): интеграция StatsCollector в receiver --- создание сессий, логирование пакетов, обработка ошибок (неверный ключ), все 5 каналов.
    \item \textbf{Команда report} (\texttt{test\_report\_cli.py}, 5 тестов): генерация отчётов, JSON-файлов и графиков через CLI.
\end{enumerate}

Общее количество тестов составляет 269 (134 модульных + 68 интеграционных + 67 тестов целостности данных). Все тесты завершаются успешно.

\subsubsection{Обработка ошибок и граничные случаи}

При разработке системы \texttt{NetStego} значительное внимание уделено корректной обработке ошибок и граничных случаев, возникающих на различных этапах конвейера обработки данных.

На этапе криптографической обработки основным источником ошибок является использование некорректного ключа при дешифровании. Алгоритм AES-256-GCM обеспечивает аутентификацию данных посредством 128-битного тега, поэтому при попытке дешифрования с неверным ключом генерируется исключение \texttt{cryptography.exceptions.InvalidTag}. Система перехватывает данное исключение и формирует информативное сообщение об ошибке, а сессия в базе статистики помечается статусом <<failed>>.

На этапе фрагментации обрабатываются следующие граничные случаи. Пустой входной файл (0~байт) обрабатывается корректно: функция \texttt{fragment} формирует единственный чанк, содержащий только SHA-256 хеш пустых данных. Файл минимального размера (1~байт) также формирует один чанк с хешем и одним байтом данных. Файл, размер которого точно кратен размеру чанка, обрабатывается без создания дополнительного пустого чанка.

На этапе сборки (\texttt{Reassembler}) обрабатываются аномалии доставки чанков. При получении чанка с некорректной контрольной суммой CRC32 чанк отбрасывается. При получении дублированного чанка (с уже полученным порядковым номером) дубликат игнорируется, а счётчик дубликатов инкрементируется. При получении чанка с порядковым номером, выходящим за пределы ожидаемого диапазона, чанк также отбрасывается.

На сетевом уровне обрабатываются ошибки привилегий, ошибки доступности интерфейса и ошибки записи PCAP-файлов. Конвейер приёма содержит пять точек возможного отказа (отсутствие пакетов, ошибка декодирования, отсутствие чанков, неполная сборка, ошибка расшифровки), и в каждой из них сессия корректно завершается со статусом <<failed>> в базе статистики.

Модуль обнаружения обрабатывает случай недостаточного объёма данных для статистического анализа. При количестве пакетов менее~10 статистические тесты возвращают результат <<Insufficient data>> без генерации ложных срабатываний. ML-детекторы (Isolation Forest и Random Forest) требуют минимум 10~признаковых векторов для построения модели; при меньшем количестве ML-анализ пропускается с соответствующим уведомлением.

\subsubsection{Анализ безопасности и ограничения системы}

Система \texttt{NetStego} разработана как исследовательский инструмент для изучения уязвимостей сетевых протоколов к стеганографическим атакам. При оценке безопасности системы необходимо учитывать как её сильные стороны, так и известные ограничения.

К сильным сторонам криптографической подсистемы относятся следующие аспекты. Использование алгоритма AES-256-GCM обеспечивает одновременно конфиденциальность и аутентичность передаваемых данных. Каждая операция шифрования использует уникальный 96-битный nonce, генерируемый криптографически стойким генератором \texttt{secrets.token\_bytes}, что исключает повторное использование nonce при разумных объёмах передачи (вероятность коллизии пренебрежимо мала при количестве операций до $2^{48}$). Функция деривации ключа Argon2id с параметрами memory\_cost~=~65536~КиБ и time\_cost~=~3 обеспечивает устойчивость к атакам перебора на GPU и ASIC.

Среди ограничений системы следует выделить несколько существенных аспектов. Во-первых, протокол NST не реализует ARQ (Automatic Repeat reQuest), что обусловлено требованием минимизации обратного трафика для повышения скрытности. Для компенсации потерь реализован механизм избыточности (параметр \texttt{-{}-redundancy}): каждый чанк дублируется $N$~раз, а приёмник дедуплицирует полученные чанки по порядковому номеру. Тестирование на канале ICMP~Payload с трёхкратной избыточностью показало восстановление данных при потере до 20\,\% пакетов. Во-вторых, сигнатура протокола NST (магическая последовательность \texttt{NST\textbackslash x00}) является обнаруживаемой при содержательном анализе полезной нагрузки. Для повышения скрытности рекомендуется использование каналов с малой ёмкостью (IP~ID, TCP~ISN, TCP~Timestamp), где данные встраиваются в заголовки и не подвергаются содержательному анализу DPI.

В-третьих, система не реализует механизмов маскировки временных характеристик трафика. Хотя джиттер частично рандомизирует межпакетные интервалы, при длительном наблюдении регулярность стеганографического трафика может быть выявлена методами анализа временных рядов. В-четвёртых, система не обеспечивает защиту от активных нарушителей (active wardens), способных модифицировать или отбрасывать подозрительные пакеты в реальном времени~[20]. Защита от активных нарушителей требует реализации механизмов помехоустойчивого кодирования и избыточности, что является направлением будущих исследований.

\subsubsection{Информационно-теоретическая ёмкость каналов}

Теоретическая пропускная способность каждого канала определяется количеством бит данных, встраиваемых в один пакет, и максимальной частотой отправки пакетов, не вызывающей подозрений. Эффективная пропускная способность $C_{\text{eff}}$ канала рассчитывается по формуле:
\begin{equation}
C_{\text{eff}} = b \cdot r \cdot (1 - \eta),
\label{eq:channel_capacity}
\end{equation}
где $b$ --- количество бит данных на пакет, $r$ --- частота отправки пакетов (PPS), $\eta$ --- доля накладных расходов протокола (заголовок NST + фрейминг). Для канала ICMP Payload при $b = 11760$~бит, $r = 10$~PPS и $\eta \approx 0{,}014$ (21~байт накладных расходов из 1470~байт) эффективная пропускная способность составляет $C_{\text{eff}} \approx 115{,}9$~Кбит/с. Для канала IP~ID при $b = 16$~бит, $r = 10$~PPS и $\eta = 0$ (данные встраиваются непосредственно в поле заголовка) $C_{\text{eff}} = 160$~бит/с $\approx 20$~байт/с.

С учётом накладных расходов протокола NST (17~байт заголовка + 4~байта фрейминга + 32~байта SHA-256 в первом чанке) и шифрования AES-256-GCM (12~байт nonce + 16~байт тега аутентификации) суммарные накладные расходы для передачи файла размером $S$~байт через канал с размером данных чанка $d$~байт составляют:
\begin{equation}
O = 28 + \lceil S_c / d \rceil \cdot 21 + 32,
\label{eq:overhead}
\end{equation}
где $S_c = S + 28$ --- размер шифротекста (включая nonce и тег), $\lceil S_c / d \rceil$ --- количество чанков, 21 = 17 + 4 --- заголовок чанка и фрейминг, 32 --- SHA-256 хеш в первом чанке.

\subsubsection{Результаты производительности}

Производительность каналов измерена путём выполнения полного конвейера обработки: шифрование, фрагментация, кодирование в пакеты, декодирование, сборка и дешифрование. Каждый тест повторялся трёхкратно для усреднения. Результаты для различных объёмов данных приведены в таблице~\ref{tab:performance}.

\begin{table}[htb]
\centering
\caption{Результаты тестирования производительности конвейера}
\label{tab:performance}
\small
\begin{tabular}{|l|c|c|c|c|}
\hline
Канал & Бит/пакет & Пакетов & Обработка & Пропуск. способн. \\
      &           & (1 КБ)  & (1 КБ)    & (Б/с) \\
\hline
ICMP Payload   & 11\,776 & 1   & 0{,}62 мс & 1\,668\,169 \\
DNS QNAME      & 240   & 44  & 15{,}8 мс & 64\,982 \\
TCP ISN        & 32    & 320 & 44{,}3 мс & 23\,116 \\
TCP Timestamp  & 32    & 320 & 63{,}2 мс & 19\,184 \\
IP ID          & 16    & 638 & 316{,}5 мс & 3\,238 \\
\hline
\end{tabular}
\end{table}

С учётом сетевых задержек при скорости отправки 10~PPS реальное время передачи файла объёмом 1~КБ составляет: ICMP~--- менее 1~с (1~пакет), DNS~--- 4{,}4~с (44~пакета), TCP~ISN и TCP~TS~--- 32~с (320~пакетов), IP~ID~--- 63{,}8~с (638~пакетов). Для файла объёмом 16~КБ количество пакетов возрастает до 12 (ICMP), 640 (DNS), 4790 (TCP~ISN, TCP~TS) и 9578 (IP~ID).

Результаты подтверждают фундаментальный компромисс между пропускной способностью и скрытностью: канал ICMP~Payload обеспечивает наивысшую пропускную способность (14\,720~Б/с при 10~PPS), однако отличается низкой скрытностью вследствие крупного объёма полезной нагрузки. Каналы IP~ID, TCP~ISN и TCP~Timestamp обеспечивают высокую скрытность за счёт встраивания в стандартные поля заголовков, однако пропускная способность ограничена 20--40~Б/с при 10~PPS.

На рисунке~\ref{fig:throughput} представлено сравнение пропускной способности каналов в логарифмической шкале.

\begin{figure}[htb]
\centering
\includegraphics[width=0.85\textwidth]{throughput_comparison.png}
\caption{Сравнение пропускной способности стеганографических каналов}
\label{fig:throughput}
\end{figure}

На рисунке~\ref{fig:timing} показана декомпозиция времени обработки по этапам конвейера для каждого канала.

\begin{figure}[htb]
\centering
\includegraphics[width=0.85\textwidth]{timing_breakdown.png}
\caption{Декомпозиция времени обработки по этапам конвейера}
\label{fig:timing}
\end{figure}

На рисунке~\ref{fig:channel_capacity} представлена ёмкость каналов (байт на пакет) в логарифмической шкале, наглядно демонстрирующая разницу в три порядка между ICMP~Payload (1470~байт) и IP~ID (2~байта).

\begin{figure}[htb]
\centering
\includegraphics[width=0.8\textwidth]{channel_capacity.png}
\caption{Сравнение ёмкости стеганографических каналов (байт/пакет)}
\label{fig:channel_capacity}
\end{figure}

На рисунке~\ref{fig:stealth_throughput} представлена визуализация компромисса между скрытностью и пропускной способностью. Каналы с высокой скрытностью (IP~ID, TCP~ISN, TCP~Timestamp) расположены в верхней части графика, но обладают низкой пропускной способностью. Канал ICMP~Payload, напротив, обеспечивает максимальную пропускную способность, но легко обнаруживается при анализе размера payload.

\begin{figure}[htb]
\centering
\includegraphics[width=0.85\textwidth]{stealth_vs_throughput.png}
\caption{Компромисс между скрытностью и пропускной способностью каналов}
\label{fig:stealth_throughput}
\end{figure}

Зависимость количества пакетов от размера передаваемого файла представлена на рисунке~\ref{fig:packets_sizes}. Для канала ICMP~Payload количество пакетов растёт медленно (файл 64~КБ --- 45~пакетов), тогда как для IP~ID рост линеен и составляет порядка 34\,000~пакетов для файла того же размера.

\begin{figure}[htb]
\centering
\includegraphics[width=0.85\textwidth]{packets_for_sizes.png}
\caption{Зависимость числа пакетов от размера файла для каждого канала}
\label{fig:packets_sizes}
\end{figure}

На рисунке~\ref{fig:overhead} показано соотношение полезных данных и накладных расходов протокола при передаче файла объёмом 1~КБ. Для канала ICMP накладные расходы составляют менее 6\,\% от общего объёма, тогда как для каналов с фиксированным размером чанка (128~байт) --- до 15\,\%.

\begin{figure}[htb]
\centering
\includegraphics[width=0.8\textwidth]{overhead_analysis.png}
\caption{Соотношение полезных данных и накладных расходов протокола (файл 1~КБ)}
\label{fig:overhead}
\end{figure}

На рисунке~\ref{fig:transfer_time} представлено время передачи файла объёмом 1~КБ при скорости 10~PPS для каждого канала.

\begin{figure}[htb]
\centering
\includegraphics[width=0.8\textwidth]{transfer_time_1kb.png}
\caption{Время передачи файла 1~КБ при скорости 10~PPS}
\label{fig:transfer_time}
\end{figure}

\subsubsection{Результаты обнаружения}

Эффективность модуля обнаружения оценена путём автоматизированного мониторинга. Для каждого канала генерировался стеганографический трафик реальным конвейером системы и нормальный трафик, имитирующий поведение ОС (инкрементальные IP~ID для Windows, монотонные TSval, реальные DNS-имена). Результаты обнаружения представлены в таблице~\ref{tab:detection_results}.

\begin{table}[htb]
\centering
\caption{Результаты тестирования модуля обнаружения (комбинированный метод)}
\label{tab:detection_results}
\begin{tabular}{|l|c|c|c|c|}
\hline
Канал & Precision & Recall & F1 & FPR \\
\hline
ICMP Payload   & $\geq$0{,}80 & $\geq$0{,}80 & $\geq$0{,}80 & $\leq$0{,}10 \\
IP ID          & $\geq$0{,}90 & $\geq$0{,}90 & $\geq$0{,}90 & $\leq$0{,}10 \\
TCP ISN        & $\geq$0{,}80 & $\geq$0{,}80 & $\geq$0{,}80 & $\leq$0{,}10 \\
DNS QNAME      & $\geq$0{,}90 & $\geq$0{,}90 & $\geq$0{,}90 & $\leq$0{,}10 \\
TCP Timestamp  & $\geq$0{,}80 & $\geq$0{,}80 & $\geq$0{,}80 & $\leq$0{,}15 \\
\hline
\end{tabular}
\end{table}

Помимо комбинированного метода, проводилась оценка отдельных подходов к обнаружению:

\begin{enumerate}[label=\arabic*)]
    \item \textbf{Статистический метод} --- наиболее эффективен для каналов ICMP (аномалия размера payload и побайтовой энтропии) и TCP~Timestamp (нарушение монотонности TSval).
    \item \textbf{ML unsupervised (Isolation Forest)} --- эффективен для обнаружения IP~ID (нарушение последовательности) и TCP~ISN (аномальное распределение ISN).
    \item \textbf{ML supervised (Random Forest)} --- при обучении на размеченных данных демонстрирует наивысшую точность классификации по всем каналам, однако требует наличия обучающей выборки.
    \item \textbf{Сигнатурный метод} --- обнаруживает DNS QNAME (base32-паттерны, терминационные маркеры) и ICMP (магическая последовательность NST в payload).
\end{enumerate}

Комбинация четырёх методов обнаружения позволяет достичь уровня ложноположительных срабатываний не более 10\,\% при обнаружении всех реализованных каналов.

На рисунке~\ref{fig:detection_accuracy} представлены результаты обнаружения в виде группированной столбчатой диаграммы Precision/Recall/F1 по каналам.

\begin{figure}[htb]
\centering
\includegraphics[width=0.85\textwidth]{detection_accuracy.png}
\caption{Точность обнаружения стеганографических каналов (Precision/Recall/F1)}
\label{fig:detection_accuracy}
\end{figure}

На рисунке~\ref{fig:redundancy} представлены результаты тестирования устойчивости к потере пакетов при трёхкратной избыточности на канале ICMP~Payload.

\begin{figure}[htb]
\centering
\includegraphics[width=0.8\textwidth]{redundancy_results.png}
\caption{Устойчивость к потере пакетов (канал ICMP, избыточность $\times$3)}
\label{fig:redundancy}
\end{figure}

\subsection{Анализ целесообразности сокрытия данных в сетевых пакетах}

Проведённое тестирование позволяет сформулировать рекомендации по выбору метода сокрытия данных в зависимости от сценария применения.

\subsubsection{Сокрытие в заголовках vs в полезной нагрузке}

Реализованные каналы делятся на два класса по способу встраивания: модификация полей заголовков (IP~ID, TCP~ISN, TCP~Timestamp) и размещение данных в полезной нагрузке пакетов (ICMP~Payload, DNS~QNAME). Каждый класс обладает принципиально различными характеристиками.

Каналы на основе заголовков (IP~ID, TCP~ISN, TCP~Timestamp) встраивают данные в поля, которые обязательно присутствуют в каждом пакете и обрабатываются стандартным сетевым стеком. Их преимущества: данные не видны системам DPI, анализирующим содержимое payload; модификация заголовков не увеличивает размер пакета; пакеты неотличимы от нормального трафика по формальным признакам (размер, протокол, порты). Недостатки: крайне низкая ёмкость (2--4~байта на пакет) и уязвимость к статистическому анализу распределения значений полей.

Каналы на основе payload (ICMP~Payload, DNS~QNAME) размещают данные в области полезной нагрузки, что обеспечивает значительно более высокую ёмкость (до 1470~байт для ICMP). Однако содержимое payload доступно для анализа системами DPI, и аномальный размер или высокая энтропия payload являются обнаруживаемыми индикаторами.

\subsubsection{Рекомендации по выбору канала}

На основе полученных результатов можно выделить следующие сценарии применения:

\begin{enumerate}[label=\arabic*)]
    \item \textbf{Передача управляющих команд C2} (объём до 100~байт): рекомендуется канал TCP~ISN или IP~ID. Высокая скрытность компенсирует низкую пропускную способность, а малый объём данных не требует значительного числа пакетов.
    \item \textbf{Передача криптографических ключей} (32--256~байт): рекомендуется канал TCP~Timestamp или TCP~ISN с избыточностью. Критичны целостность и скрытность, а не скорость передачи.
    \item \textbf{Эксфильтрация документов} (1--100~КБ): рекомендуется канал DNS~QNAME --- баланс между ёмкостью (30~байт/пакет) и маскировкой под легитимный DNS-трафик. DNS практически не блокируется в корпоративных сетях.
    \item \textbf{Передача больших файлов} ({>}100~КБ): рекомендуется канал ICMP~Payload при отсутствии активного мониторинга или с маскировочным ICMP-трафиком.
\end{enumerate}

Следует отметить, что целесообразность использования конкретного канала определяется не только его техническими характеристиками, но и особенностями целевой сетевой среды: наличием межсетевых экранов, блокирующих ICMP; использованием NAT, модифицирующего IP~ID; применением DNS-прокси, кеширующих запросы. Система \texttt{NetStego} позволяет исследовать каждый из этих факторов путём анализа PCAP-дампов реального трафика.

\subsection{Демонстрация работы системы}

Для подтверждения работоспособности системы проведена практическая демонстрация полного цикла передачи скрытых данных: генерация ключа, отправка сообщения, приём и расшифровка, обнаружение скрытого канала.

\subsubsection{Генерация ключа шифрования}

На первом этапе генерируется 256-битный ключ шифрования AES-256-GCM с помощью команды \texttt{keygen}. Ключ сохраняется в бинарный файл \texttt{demo\_key.bin} и используется на всех последующих этапах передачи и приёма данных (рисунок~\ref{fig:demo_keygen}).

\begin{figure}[htb]
\centering
\includegraphics[width=0.85\textwidth]{demo_keygen.png}
\caption{Генерация 256-битного ключа шифрования}
\label{fig:demo_keygen}
\end{figure}

\subsubsection{Отправка сообщения через канал ICMP Payload}

Текстовое сообщение <<Secret message for NetStego demo>> (32~байта) передаётся через канал ICMP~Payload. Система выполняет шифрование AES-256-GCM (результат --- 60~байт шифротекста), фрагментацию на чанки с заголовком NST и кодирование в ICMP Echo Request пакет. Сформированные пакеты сохраняются в файл \texttt{demo\_icmp.pcap} (рисунок~\ref{fig:demo_send_icmp}).

\begin{figure}[htb]
\centering
\includegraphics[width=0.85\textwidth]{demo_send_icmp.png}
\caption{Отправка текстового сообщения через канал ICMP Payload}
\label{fig:demo_send_icmp}
\end{figure}

\subsubsection{Приём и расшифровка сообщения (ICMP)}

Приём данных выполняется из сохранённого PCAP-файла с помощью команды \texttt{receive}. Система последовательно выполняет: чтение пакетов из PCAP, декодирование каналом ICMP, извлечение и сборку чанков, верификацию CRC32 и SHA-256, расшифрование AES-256-GCM. На рисунке~\ref{fig:demo_receive_icmp} видно, что из 1~пакета декодировано 113~байт, извлечён 1~чанк, собрано 60~байт шифротекста, после расшифровки получено 32~байта исходного текста. Команда \texttt{type received.txt} подтверждает побайтовое совпадение с исходным сообщением: <<Secret message for NetStego demo>>.

\begin{figure}[htb]
\centering
\includegraphics[width=0.85\textwidth]{demo_receive_icmp.png}
\caption{Приём и расшифровка сообщения из канала ICMP --- подтверждение целостности}
\label{fig:demo_receive_icmp}
\end{figure}

\subsubsection{Отправка файла через канал DNS QNAME}

Для демонстрации работы канала DNS QNAME в качестве входных данных используется текстовый файл. Система кодирует зашифрованные данные в base32-поддомены DNS-запросов. Ввиду меньшей ёмкости канала (30~байт на пакет) данные разбиваются на 3~чанка и формируют соответствующее количество DNS-пакетов (рисунок~\ref{fig:demo_send_dns}).

\begin{figure}[htb]
\centering
\includegraphics[width=0.85\textwidth]{demo_send_dns.png}
\caption{Отправка файла через канал DNS QNAME}
\label{fig:demo_send_dns}
\end{figure}

\subsubsection{Обнаружение скрытого канала в ICMP-трафике}

Модуль обнаружения анализирует файл \texttt{demo\_icmp.pcap} и выносит вердикт <<SUSPICIOUS>> с уровнем уверенности 70\,\%. Статистический анализ выявляет аномальный размер ICMP payload (115~байт --- нехарактерно для стандартных ICMP Echo Request). Сигнатурный анализ обнаруживает магическую последовательность \texttt{NST\textbackslash x00} протокола фрагментации в пакете~\#0 на смещении~6 (рисунок~\ref{fig:demo_detect_icmp}).

\begin{figure}[htb]
\centering
\includegraphics[width=0.85\textwidth]{demo_detect_icmp.png}
\caption{Результат обнаружения скрытого канала в ICMP-трафике}
\label{fig:demo_detect_icmp}
\end{figure}

\subsubsection{Обнаружение скрытого канала в DNS-трафике}

Анализ файла \texttt{demo\_dns.pcap} также выносит вердикт <<SUSPICIOUS>> с уровнем уверенности 70\,\%. Сигнатурный метод обнаруживает характерные для стеганографии base32-паттерны в DNS QNAME и терминационные маркеры. Статистический метод выявляет высокую энтропию поддоменов, нехарактерную для легитимного DNS-трафика (рисунок~\ref{fig:demo_detect_dns}).

\begin{figure}[htb]
\centering
\includegraphics[width=0.85\textwidth]{demo_detect_dns.png}
\caption{Результат обнаружения скрытого канала в DNS-трафике}
\label{fig:demo_detect_dns}
\end{figure}

\subsubsection{Автоматизированное тестирование}

Для подтверждения корректности работы всех компонентов системы выполнен запуск полного набора автоматизированных тестов (рисунок~\ref{fig:demo_pytest}). Результат --- 269~тестов пройдено успешно (134~модульных, 68~интеграционных, 67~тестов целостности данных). Тесты покрывают все стеганографические каналы, криптографический модуль, протокол фрагментации и сборки, модуль обнаружения, систему бенчмарков и визуализации.

\begin{figure}[htb]
\centering
\includegraphics[width=0.85\textwidth]{demo_pytest.png}
\caption{Результат запуска автоматизированных тестов (269 passed)}
\label{fig:demo_pytest}
\end{figure}

Таким образом, демонстрация подтверждает полную работоспособность конвейера: исходное сообщение проходит этапы шифрования, фрагментации, стеганографического кодирования и обратного преобразования с побайтовым совпадением результата. Модуль обнаружения успешно выявляет скрытые каналы как в ICMP, так и в DNS-трафике.

\subsection{Выводы по реализации}

В результате работы реализована модульная и расширяемая система \texttt{NetStego}, включающая:
\begin{enumerate}[label=\arabic*)]
    \item пять стеганографических каналов с различными характеристиками ёмкости и скрытности, включая конфигурируемый DNS-суффикс для снижения обнаруживаемости;
    \item криптографическую защиту на основе AES-256-GCM с деривацией ключа Argon2id;
    \item протокол фрагментации с контролем целостности (CRC32, SHA-256) и поддержкой out-of-order сборки;
    \item многоуровневый модуль обнаружения, объединяющий статистический анализ (9~тестов), ML-обнаружение (Isolation Forest и Random Forest с 9~признаками) и сигнатурный поиск;
    \item модуль анализа содержимого payload для обнаружения вредоносных вложений (20~сигнатур в 7~категориях);
    \item систему сбора статистики с базой данных SQLite, интегрированную в конвейеры отправки и приёма;
    \item автоматическую генерацию 6~типов графиков и отчётов для демонстрации результатов;
    \item удобный CLI-интерфейс с 7~командами и поддержкой экспорта в JSON, CSV и HTML.
\end{enumerate}

Тестирование подтвердило корректность работы всех компонентов (269~тестов) и позволило оценить производительность каналов и эффективность методов обнаружения.
